% =====================================================================
% Standalone 4-page project summary (EN pages 1-2, AR pages 3-4).
% Matches the main report's class, fonts, and general look.
% Compile with XeLaTeX:  xelatex summary.tex
% Content source: report/two_page_summary.md (verified numbers).
% =====================================================================
\documentclass[11pt,a4paper]{report}

\usepackage{graphicx}
\usepackage{amsmath}
\usepackage[left=1.9cm,top=1.7cm,right=1.9cm,bottom=1.7cm]{geometry}
\usepackage{newtxmath}
\usepackage{fontspec}
\usepackage{polyglossia}
\setmainlanguage{english}
\setotherlanguage{arabic}
\newfontfamily\arabicfont[Script=Arabic,Scale=1.1]{Traditional Arabic}
\usepackage{microtype}
\usepackage[hidelinks]{hyperref}
\linespread{0.96}
\setlength{\parindent}{0pt}
\setlength{\parskip}{2pt}
\pagestyle{plain}

% Compact numbered run-in heading for the summary.
\newcommand{\shead}[1]{\par\vspace{3pt}\noindent\textbf{#1}\par\vspace{1pt}}

\begin{document}

% ============================ ENGLISH, PAGES 1-2 =====================
\enlargethispage{3\baselineskip}
\begin{center}
{\large \textbf{Voice Signal Analysis Using Machine Learning for Early
Detection of Parkinson's Disease}}\\[2pt]
{\small Suhail Mohamed Alhammali --- ID 2200208522 ---
Supervisor: Dr.\ Adel Adyaf\\
Control and Instrumentation Division, Department of Biomedical
Engineering,\\
Faculty of Engineering, University of Tripoli --- Spring 2026}
\end{center}

\shead{1. Motivation}
Parkinson's disease is the second most common neurodegenerative
disorder, affecting an estimated 6.1 million people worldwide in 2016,
more than double the 1990 figure. It is caused by the progressive loss
of dopamine-producing neurons in brain regions that control movement,
and by the time the characteristic motor signs appear, a substantial
proportion of those neurons has already degenerated. No routine
laboratory test exists; diagnosis remains clinical.

The same impairment that affects the limbs also affects the small
muscles of respiration, the larynx, and articulation. The resulting
speech disorder, hypokinetic dysarthria, is present in roughly 90\% of
patients, and quantitative acoustic analysis has detected vocal
impairment in about 78\% of patients with early, untreated disease. A
longitudinal case study reported measurable reduction of pitch
variability in recordings made around five years before clinical
diagnosis.

A microphone is inexpensive, non-invasive, and repeatable. If acoustic
measurements can be related to disease status with useful reliability,
voice analysis could act as a low-cost screening-support signal,
flagging individuals for whom a clinical examination is advisable.

\shead{2. Objectives}
The project set out to inspect and verify a public voice corpus before
any modelling; to implement, as a core engineering contribution, the
extraction of acoustic features directly from raw audio; to train and
compare classical, explainable classifiers; to evaluate them under
strictly subject-independent validation; to measure rather than assume
generalisation by testing on a second, foreign corpus; and to deliver a
prototype whose output communicates its own uncertainty.

\shead{3. Methodology}
\textbf{Data.} The MDVR-KCL corpus provides 73 smartphone recordings
from 37 English-speaking participants --- 21 healthy controls and 16
patients --- each performing a read-text task and a
spontaneous-dialogue task. All files are mono WAV at 44.1~kHz, between
73 and 221 seconds long. Before any modelling, an automated inspection
verified file integrity, absence of duplicate audio by MD5 comparison,
agreement between the two independent sources of each class label, and
the presence of a parseable speaker identifier in every filename. All
73 files passed.

\textbf{Preprocessing.} Each recording is conditioned in five fixed
steps: mono conversion, resampling to 44.1~kHz, DC-offset removal,
trimming of leading and trailing silence 30~dB below peak, and peak
normalisation to 0.95. No denoising is applied. This is a deliberate
scientific decision: denoising algorithms work by detecting and
suppressing signal irregularity, and cycle-to-cycle irregularity is
precisely the clinical information that the perturbation measures
quantify. Aggressive cleaning would make a disordered voice appear
artificially healthy.

\textbf{Features.} Each recording is divided into consecutive
10-second chunks, yielding 1001 chunks in total, and 74 acoustic
features are extracted from every chunk: fundamental-frequency
statistics (7), jitter variants (4), shimmer variants (5),
harmonics-to-noise ratio (1), smoothed cepstral peak prominence (1),
pause and timing statistics (4), and MFCC and delta-MFCC summaries
(52). Each family is chosen for a physiological reason --- pitch
variability for monotone speech, jitter and shimmer for unstable
vocal-fold control, harmonics-to-noise ratio and cepstral peak
prominence for incomplete glottal closure, pause statistics for
disturbed speech timing, and spectral features for imprecise
articulation. The recording is represented by the mean of its chunks'
feature vectors. All extraction was implemented within this project; no
pre-extracted feature set was used. One implementation serves training,
evaluation, and prediction alike, enforced by a saved configuration
contract verified before every prediction.

\textbf{Validation.} Every classifier runs inside a pipeline of median
imputation, standardisation, and the model itself, so that all
statistics learned from data are fitted on training folds only.
Evaluation uses stratified grouped 5-fold cross-validation with the
speaker identity as the grouping variable, repeated with three random
seeds, so that no person's recordings ever appear on both sides of a
split. A runtime assertion recomputes the training/test speaker
intersection in every fold and halts the program if it is non-empty; it
never triggered. Balanced accuracy is the primary metric, because a
model that always answers ``healthy'' would score 0.575 plain accuracy
while detecting no patient at all.

\textbf{Selection rules, declared before experimentation.} Identical
splits for every variant; subject-level balanced accuracy as the
primary metric; a more complex variant adopted only if it exceeds the
simpler one by more than 0.02; any result at or above 0.95 halting the
study for a leakage investigation; and full reporting of all
configurations, including rejected ones.

\enlargethispage{4\baselineskip}%
\shead{4. Results}
Nineteen configurations were evaluated. The adopted model is a
default-settings Random Forest on the 74 chunk-averaged features,
reaching a subject-level balanced accuracy of $0.822 \pm 0.032$,
sensitivity 0.771, specificity 0.873, and ROC-AUC 0.864 --- an
improvement of 0.042 over the whole-recording baseline. A soft-voting
ensemble scored higher at 0.840, but was rejected because it exceeded
the winner by only 0.018, less than the pre-declared margin and smaller
than the observed seed-to-seed variability of about 0.03.
Hyperparameter tuning did not improve on default settings at all; the
genuine gains came from the continuous-speech features and chunk
averaging.

The model is interpretable. Its most important feature is the
fundamental-frequency range, the direct acoustic correlate of the
monotone, prosodically flattened speech that clinical descriptions
identify as a hallmark of the disease, followed by the pitch maximum
and by variability measures of the spectral envelope and its dynamics.
Variability features consistently outrank averages, which is consistent
with reduced articulatory movement rather than a shifted average
spectrum.

A confounder check was also performed. Because the corpus carries no
verified sex metadata and absolute pitch partly encodes speaker sex,
the whole evaluation was repeated with the four absolute-pitch features
removed. The selected model was unaffected (0.768 to 0.780), showing
that its decisions rest on perturbation and spectral information; a
small neural network, by contrast, collapsed from 0.671 to 0.526,
revealing that it had been relying on the potential confound, and was
excluded from further work.

\shead{5. Two findings that carry more weight than the headline number}
\textbf{Validation design dominates reported performance.} The
identical pipeline and model were evaluated twice, changing only the
split: 0.775 balanced accuracy under the correct speaker-grouped split
against 0.909 under a deliberately incorrect chunk-level split. The
0.13 inflation is produced entirely by the split --- it measures
recognition of recordings the model has partly seen, not detection of
disease --- and it exceeds all genuine improvements achieved in this
project combined. This is the measured explanation for why many
published results on comparable data report 95\% or more.

\textbf{Generalisation was measured, not assumed.} The frozen model,
with no retraining and no threshold adjustment, was tested on an
independent Italian corpus of 61 speakers in a different language,
recorded with different equipment. Inspection first revealed a trap:
every 44.1~kHz file in that corpus belongs to the patient group while
all healthy files are 16~kHz, so a model could have separated the
groups by recording bandwidth alone. All audio was therefore
band-limited to 16~kHz before evaluation. The result was 0.701 ROC-AUC
and 0.629 balanced accuracy on the age-fair comparison --- evidence
that a genuine, transferable acoustic signal survives domain shift,
together with an honest measurement of how much is lost to it.

\shead{6. Prototype}
A browser application and a command-line tool share a single prediction
function, so the two interfaces cannot diverge. Scores between 0.35 and
0.65 are reported as inconclusive rather than as a class; warnings are
attached for sample-rate mismatch, clipping, low signal-to-noise ratio,
and very short recordings; and no audio is retained, uploaded files
being deleted immediately after analysis. In the external test, this
design placed 100\% of elderly healthy speakers in the inconclusive
band instead of confidently mislabelling them --- the intended
behaviour for input far outside the system's experience.

\shead{7. Limitations and scope}
The dominant limitation is sample size: 37 training subjects from one
recording setup and one language, without demographic metadata or
disease severity, so the system's sensitivity to early cases
specifically remains unmeasured. The perturbation measures are used on
continuous speech although they were defined for sustained vowels.
Scores are not calibrated probabilities. One external corpus is a
single data point about generalisation, not a characterisation of it.

The system is a non-diagnostic research screening-support prototype. It
reports the similarity of an acoustic pattern to dataset groups. It
does not diagnose Parkinson's disease, does not measure any person's
medical risk, and cannot replace evaluation by a qualified healthcare
professional.

\shead{8. Conclusion}
On small voice corpora, the choice of validation scheme moves reported
performance more than model engineering does. An honestly validated
0.822 on unseen people, accompanied by a measured cross-corpus 0.701
AUC and an interface that declines to answer when uncertain, is a
stronger scientific result than an inflated figure --- and the
appropriate foundation for any future clinical development.

\clearpage

% ============================ ARABIC, PAGES 3-4 ======================
\begin{Arabic}

\begin{center}
{\large \textbf{تحليل إشارة الصوت باستخدام تعلّم الآلة للكشف المبكر عن مرض باركنسون}}

\vspace{3pt}
سهيل محمد الحمالي، رقم القيد \textenglish{2200208522}

المشرف: الدكتور عادل أضياف

شعبة التحكم والقياس، قسم الهندسة الطبية الحيوية، كلية الهندسة، جامعة طرابلس

فصل ربيع \textenglish{2026}
\end{center}
\vspace{2pt}

\shead{1. الدافع}
مرض باركنسون هو ثاني أكثر الأمراض التنكّسية العصبية شيوعاً، وقُدِّر عدد
المصابين به في العالم بنحو \textenglish{6.1} مليون شخص سنة 2016، أي
أكثر من ضعف عددهم سنة 1990. وسببه الفقدان التدريجي للخلايا العصبية
المنتجة للدوبامين في مناطق الدماغ المسؤولة عن التحكم بالحركة. وعند ظهور
العلامات الحركية المميّزة يكون جزء كبير من تلك الخلايا قد فُقد بالفعل.
ولا يوجد فحص مخبري روتيني للمرض، فالتشخيص يبقى سريرياً.

والخلل نفسه الذي يصيب الأطراف يصيب أيضاً العضلات الصغيرة المسؤولة عن
التنفّس والحنجرة ومخارج الحروف. ويُسمّى اضطراب الكلام الناتج عن ذلك
«عُسر التلفّظ ناقص الحركة»، وهو موجود لدى نحو تسعين بالمئة من المرضى.
وقد كشف التحليل الصوتي الكمّي وجود اضطراب صوتي لدى نحو ثمانية وسبعين
بالمئة من المرضى في المرحلة المبكّرة غير المعالَجة. كما أشارت دراسة
حالة طولية إلى انخفاض قابل للقياس في تغيّرية النبرة في تسجيلات أُخذت
قبل التشخيص السريري بنحو خمس سنوات.

والميكروفون رخيص وغير جارح ويمكن تكراره بسهولة. فإذا أمكن ربط القياسات
الصوتية بحالة المرض بموثوقية مفيدة، فقد يعمل تحليل الصوت إشارةً لدعم
الفرز منخفض التكلفة، تُشير إلى الأشخاص الذين يُستحسن أن يخضعوا لفحص
سريري.

\shead{2. الأهداف}
هدف المشروع إلى فحص قاعدة بيانات صوتية عامة والتحقّق منها قبل أي نمذجة،
وإلى تنفيذ استخراج الخصائص الصوتية من الصوت الخام بوصفه المساهمة
الهندسية الأساسية، وإلى تدريب مصنّفات كلاسيكية قابلة للتفسير
ومقارنتها، وإلى تقييمها بتحقّق مستقل تماماً عن المتحدث، وإلى قياس قدرة
التعميم بدل افتراضها عبر الاختبار على قاعدة بيانات أجنبية ثانية، وإلى
تقديم نموذج أولي يُبلّغ عن عدم يقينه.

\shead{3. المنهجية}
\textbf{البيانات.} توفّر قاعدة البيانات \textenglish{MDVR-KCL} ثلاثة
وسبعين تسجيلاً بالهاتف المحمول من سبعة وثلاثين مشاركاً ناطقاً
بالإنجليزية: واحد وعشرون سليماً وستة عشر مريضاً، أدّى كل منهم مهمة
قراءة نصّ ومهمة حوار عفوي. وجميع الملفات بصيغة \textenglish{WAV}
أحادية القناة بمعدّل \textenglish{44.1} كيلوهرتز، وتتراوح مدّتها بين
ثلاث وسبعين ومئتين وإحدى وعشرين ثانية. وقبل أي نمذجة، تحقّق فحص آلي من
سلامة الملفات، ومن خلوّها من التكرار عبر مقارنة بصمات
\textenglish{MD5}، ومن تطابق مصدرَي التسمية المستقلّين لكل ملف، ومن
وجود معرّف متحدث قابل للاستخراج في كل اسم ملف. واجتازت الملفات الثلاثة
والسبعون جميع الفحوص.

\textbf{المعالجة المسبقة.} يمرّ كل تسجيل بخمس خطوات ثابتة: التحويل إلى
قناة واحدة، وتوحيد معدّل العيّنات عند \textenglish{44.1} كيلوهرتز،
وإزالة الانزياح المستمر، وقصّ الصمت من بداية التسجيل ونهايته عند
ثلاثين ديسيبل تحت القمة، والتطبيع إلى قمة قدرها \textenglish{0.95}.
ولم تُطبَّق أي إزالة للضجيج، وهذا قرار علمي مقصود: فخوارزميات إزالة
الضجيج تعمل بكشف عدم الانتظام في الإشارة وكبته، وعدم الانتظام بين دورة
وأخرى هو بالضبط المعلومة السريرية التي تقيسها خصائص الاضطراب. ولو
طُبِّق تنظيف عدواني لبدا الصوت المضطرب سليماً بشكل مصطنع.

\textbf{الخصائص.} يُقطَّع كل تسجيل إلى مقاطع متتابعة مدّتها عشر ثوانٍ،
فينتج عن ذلك ألف ومقطع واحد، وتُستخلص من كل مقطع أربع وسبعون خاصية
صوتية: إحصاءات التردد الأساسي وعددها سبع، وأنواع الارتجاف وعددها
أربعة، وأنواع التذبذب وعددها خمسة، ونسبة التوافقيات إلى الضجيج، وبروز
القمة الكبسترالية المُنعَّم، وإحصاءات التوقّفات وعددها أربع، وملخّصات
معاملات الكبستروم على مقياس ميل ومشتقاتها وعددها اثنتان وخمسون. ولكل
عائلة مبرّر فسيولوجي: تغيّرية النبرة لرتابة الكلام، والارتجاف والتذبذب
لعدم استقرار التحكم في الحبال الصوتية، ونسبة التوافقيات إلى الضجيج
وبروز القمة الكبسترالية لعدم اكتمال إغلاق المزمار، وإحصاءات التوقّفات
لاضطراب إيقاع الكلام، والخصائص الطيفية لعدم دقة مخارج الحروف. ويُمثَّل
التسجيل بمتوسط متّجهات خصائص مقاطعه. وقد نُفِّذ الاستخراج كلّه ضمن هذا
المشروع، ولم تُستخدم أي مجموعة خصائص جاهزة. ويُستعمل تطبيق برمجي واحد
للتدريب والتقييم والتنبّؤ معاً، ويُفرَض ذلك بعقد إعدادات محفوظ
يُتحقَّق منه قبل كل تنبّؤ.

\textbf{التحقّق.} يعمل كل مصنّف داخل أنبوب مكوّن من سدّ القيم الناقصة
بالوسيط، ثم التقييس، ثم النموذج نفسه، بحيث تُقدَّر كل الإحصاءات
المتعلّمة من بيانات التدريب وحدها في كل طيّة. ويستخدم التقييم تحقّقاً
متقاطعاً مجمَّعاً ومطبَّقاً على خمس طيّات، بمعرّف المتحدث بوصفه متغيّر
التجميع، مكرَّراً بثلاث بذور عشوائية، بحيث لا تظهر تسجيلات أي شخص على
جانبي أي تقسيم. ويعيد تأكيد برمجي حساب تقاطع متحدّثي التدريب والاختبار
في كل طيّة، ويوقف البرنامج إن لم يكن التقاطع فارغاً، ولم يُفعَّل هذا
التأكيد قطّ. والدقة المتوازنة هي المقياس الأساسي، لأن نموذجاً يجيب
دائماً بأن الشخص سليم يحقّق دقة عادية قدرها \textenglish{0.575} دون أن
يكتشف أي مريض.

\textbf{قواعد الاختيار المُعلنة قبل التجارب.} تقسيمات متطابقة لكل
المتغيّرات، والدقة المتوازنة على مستوى الشخص مقياساً أساسياً، وعدم
اعتماد أي متغيّر أعقد إلا إذا تجاوز الأبسط بأكثر من
\textenglish{0.02}، وإيقاف الدراسة للتحقيق في تسريب البيانات عند بلوغ
أي نتيجة \textenglish{0.95} فأكثر، والإبلاغ عن كل التكوينات بما فيها
المرفوضة.

\shead{4. النتائج}
جرى تقييم تسعة عشر تكويناً. والنموذج المعتمد هو غابة عشوائية بإعدادات
افتراضية على الخصائص الأربع والسبعين المأخوذة بالمتوسط، وقد بلغ دقة
متوازنة على مستوى الشخص قدرها \textenglish{0.822} بانحراف معياري
\textenglish{0.032}، وحساسية \textenglish{0.771}، ونوعية
\textenglish{0.873}، ومساحة تحت المنحنى قدرها \textenglish{0.864}، أي
بتحسّن قدره \textenglish{0.042} عن الخط الأساسي المحسوب على التسجيل
كاملاً. وحقّق نموذج تجميعي نتيجة أعلى قدرها \textenglish{0.840}، لكنه
رُفض لأنه تجاوز الفائز بمقدار \textenglish{0.018} فقط، وهو أقل من
الهامش المُعلن مسبقاً وأصغر من التباين المرصود بين البذور والبالغ نحو
\textenglish{0.03}. ولم يُحسِّن ضبط المعاملات الفائقة النتيجة عن
الإعدادات الافتراضية إطلاقاً، بل جاء التحسّن الحقيقي من خصائص الكلام
المتّصل ومن أخذ متوسط المقاطع.

والنموذج قابل للتفسير. فأهم خصائصه هي مدى التردد الأساسي، وهو المقابل
الصوتي المباشر لرتابة النبرة وانبساط التنغيم اللذين تصفهما المراجع
السريرية بوصفهما علامة مميّزة للمرض، يليه أعلى قيمة للنبرة، ثم مقاييس
تغيّرية الغلاف الطيفي وديناميكياته. وتتفوّق خصائص التغيّرية على خصائص
المتوسط باستمرار، وهو ما يتّسق مع نقص سعة حركة أعضاء النطق لا مع
انزياح في متوسط الطيف.

وأُجري كذلك فحص للعامل المُربك. فلأن قاعدة البيانات لا تحمل بيانات
موثّقة عن جنس المتحدث، ولأن النبرة المطلقة تعكس جنس المتحدث جزئياً،
أُعيد التقييم كاملاً بعد حذف خصائص النبرة المطلقة الأربع. ولم يتأثّر
النموذج المختار، إذ انتقل من \textenglish{0.768} إلى
\textenglish{0.780}، ممّا يدلّ على أن قراراته تستند إلى معلومات
الاضطراب والطيف. أما شبكة عصبية صغيرة فقد انهارت من
\textenglish{0.671} إلى \textenglish{0.526}، ممّا كشف اعتمادها على
العامل المُربك المحتمل، فاستُبعدت من الأعمال اللاحقة.

\shead{5. نتيجتان أثقل وزناً من الرقم الرئيسي}
\textbf{تصميم التحقّق يحكم الأداء المُعلن.} قُيِّم الأنبوب نفسه
والنموذج نفسه مرتين، ولم يتغيّر سوى طريقة التقسيم: دقة متوازنة قدرها
\textenglish{0.775} بالتقسيم الصحيح المجمَّع حسب المتحدث، مقابل
\textenglish{0.909} بتقسيم خاطئ متعمَّد على مستوى المقاطع. والتضخّم
البالغ \textenglish{0.13} ناتج عن التقسيم وحده، إذ يقيس قدرة النموذج
على التعرّف على تسجيلات رآها جزئياً لا على كشف المرض، وهو يفوق مجموع
كل التحسينات الحقيقية التي تحقّقت في هذا المشروع. وهذا هو التفسير
المقيس لسبب إعلان كثير من الأبحاث المنشورة على بيانات مماثلة نتائج
تبلغ خمسة وتسعين بالمئة أو أكثر.

\textbf{قدرة التعميم قيست ولم تُفترَض.} اختُبر النموذج مجمَّداً، بلا
إعادة تدريب وبلا تعديل للعتبة، على قاعدة بيانات إيطالية مستقلة تضم
واحداً وستين متحدثاً بلغة أخرى وبأجهزة تسجيل مختلفة. وكشف الفحص أولاً
عن فخّ: كل ملف بمعدّل \textenglish{44.1} كيلوهرتز في تلك القاعدة يخصّ
مجموعة المرضى، بينما كل ملفات الأصحّاء بمعدّل \textenglish{16}
كيلوهرتز، فكان بإمكان النموذج أن يفصل بين المجموعتين بعرض النطاق
الترددي وحده. ولذلك جرى تحديد نطاق كل الملفات الصوتية عند
\textenglish{16} كيلوهرتز قبل التقييم. وكانت النتيجة مساحة تحت المنحنى
قدرها \textenglish{0.701} ودقة متوازنة قدرها \textenglish{0.629} في
المقارنة العادلة عمرياً، وهو دليل على بقاء إشارة صوتية حقيقية قابلة
للانتقال رغم انزياح المجال، مع قياس أمين لمقدار ما يُفقد بسببه.

\shead{6. النموذج الأولي}
يشترك تطبيق المتصفّح وواجهة سطر الأوامر في دالّة تنبّؤ واحدة، فلا يمكن
للواجهتين أن تختلفا. وتُعرَض النتائج الواقعة بين \textenglish{0.35}
و\textenglish{0.65} بوصفها غير حاسمة بدلاً من إعطائها تصنيفاً. وتُرفق
تحذيرات عند اختلاف معدّل العيّنات، وعند حدوث قَصّ في الإشارة، وعند
انخفاض نسبة الإشارة إلى الضجيج، وعند قِصَر التسجيل. ولا يُحتفَظ بأي
ملف صوتي، إذ تُحذف الملفات المرفوعة فور تحليلها. وفي الاختبار الخارجي
وضع هذا التصميم مئة بالمئة من المتحدثين الأصحّاء كبار السن في النطاق
غير الحاسم بدلاً من تصنيفهم تصنيفاً خاطئاً بثقة، وهو السلوك المقصود مع
مُدخَل بعيد إلى هذا الحدّ عن خبرة النظام.

\shead{7. القيود والنطاق}
القيد الأبرز هو حجم العيّنة: سبعة وثلاثون شخصاً في التدريب من إعداد
تسجيل واحد ولغة واحدة، بلا بيانات ديموغرافية وبلا معلومات عن شدّة
المرض، ولذلك تبقى حساسية النظام للحالات المبكّرة تحديداً غير مقيسة.
كما تُستخدم مقاييس الاضطراب على الكلام المتّصل رغم أنها عُرِّفت أصلاً
للحروف الممدودة. ودرجات النموذج ليست احتمالات معايَرة. وقاعدة بيانات
خارجية واحدة تمثّل نقطة بيانات واحدة عن التعميم لا وصفاً كاملاً له.

والنظام أداة أولية لدعم الفرز البحثي وغير تشخيصية. فهو يُبلّغ عن مدى
تشابه نمط صوتي مع مجموعات قاعدة البيانات، ولا يشخّص مرض باركنسون، ولا
يقيس خطراً طبياً على أي شخص، ولا يمكن أن يحلّ محل تقييم مختصّ رعاية
صحية مؤهّل.

\shead{8. الخلاصة}
على قواعد البيانات الصوتية الصغيرة، يُحرّك اختيار مخطّط التحقّق
الأداءَ المُعلن أكثر ممّا تُحرّكه هندسة النموذج. ونتيجة قدرها
\textenglish{0.822} على أشخاص لم يرهم النموذج، متحقَّق منها بأمانة،
ومصحوبة بمساحة تحت المنحنى مقيسة عبر قاعدتَي بيانات قدرها
\textenglish{0.701}، وبواجهة تمتنع عن الإجابة عند عدم اليقين، هي
نتيجة علمية أقوى من رقم مضخَّم، وهي الأساس المناسب لأي تطوير سريري
مستقبلي.

\end{Arabic}
\end{document}
